{\small
\setlength{\tabcolsep}{5pt}
\begin{longtable}{@{}>{\raggedright\arraybackslash}p{0.48\linewidth}>{\raggedright\arraybackslash}p{0.48\linewidth}@{}}
\toprule
% DOCX table after paragraph 1168, row 0, column 0
\textbf{Ondas electromagnéticas}
&
% DOCX table after paragraph 1168, row 0, column 1
\textbf{Ondas gravitatorias}
\\[6pt]
\midrule
% DOCX table after paragraph 1168, row 1, column 0
Son oscilaciones de los campos eléctrico y magnético que se propagan por el espaciotiempo.
&
% DOCX table after paragraph 1168, row 1, column 1
Son oscilaciones del propio espaciotiempo.
\\[6pt]
% DOCX table after paragraph 1168, row 2, column 0
Las más simples corresponden a la radiación\index{dipolo!radiación} dipolar, generada por un dipolo\index{dipolo!eléctrico} oscilante: dos cargas eléctricas opuestas.
&
% DOCX table after paragraph 1168, row 2, column 1
Las más simples corresponden a la radiación\index{cuadrupolo!radiación} cuadrupolar, generada por un cuadrupolo de masa\index{cuadrupolo!masa} oscilante: una distribución no circular de masa.
\\[6pt]
% DOCX table after paragraph 1168, row 3, column 0
Son casi siempre el resultado de la superposición incoherente de contribuciones de millones de electrones, átomos o moléculas.
&
% DOCX table after paragraph 1168, row 3, column 1
Se producen por el movimiento coherente de grandes cantidades de masa y energía, ya sea en forma de objetos astronómicos o de concentraciones de energía.
\\[6pt]
% DOCX table after paragraph 1168, row 4, column 0
Nos muestran los detalles del estado termodinámico\index{estado termodinámico}, temperatura y densidad, de las fuentes que las producen.
&
% DOCX table after paragraph 1168, row 4, column 1
Nos muestran los detalles de la dinámica y las propiedades globales, tamaño, masa y velocidad, de grandes concentraciones de masa-energía.
\\[6pt]
% DOCX table after paragraph 1168, row 5, column 0
Sus longitudes de onda son muy cortas en comparación con el tamaño de las fuentes astronómicas. Por tanto, es posible obtener «imágenes» de las fuentes.
&
% DOCX table after paragraph 1168, row 5, column 1
Sus longitudes de onda suelen ser comparables al tamaño de las fuentes que las producen, o incluso mayores. Por tanto, no es posible obtener «imágenes» de las fuentes.
\\[6pt]
% DOCX table after paragraph 1168, row 6, column 0
Sufren absorción, difusión y dispersión por la materia que encuentran en su camino.
&
% DOCX table after paragraph 1168, row 6, column 1
Se propagan sin perturbaciones, a la velocidad de la luz, a través de todo tipo de materiales.
\\[6pt]
% DOCX table after paragraph 1168, row 7, column 0
Las emitidas por objetos celestes tienen un espectro\index{espectro!electromagnético} esencialmente independiente de la masa de la fuente; por ejemplo, las emisiones de agujeros negros estelares y supermasivos son muy parecidas.
&
% DOCX table after paragraph 1168, row 7, column 1
Las emitidas por objetos celestes tienen un espectro que depende de la masa de la fuente: la frecuencia de las ondas gravitatorias\index{espectro!ondas gravitatorias} de agujeros negros estelares es mucho mayor que la de las procedentes de agujeros negros supermasivos.
\\[6pt]
% DOCX table after paragraph 1168, row 8, column 0
Sus longitudes de onda son menores que los radios de curvatura de los objetos que encuentran al propagarse. Por eso experimentan el efecto de lente gravitatoria\index{gravitatorio!efecto de lente}, es decir, distorsión y amplificación.
&
% DOCX table after paragraph 1168, row 8, column 1
Sus longitudes de onda son comparables o superiores a los radios de curvatura de los objetos que encuentran al propagarse. Por eso el efecto de lente gravitatoria es menos pronunciado y más difícil de medir.
\\[6pt]
\bottomrule
\end{longtable}
}
